Paper~5~\cite{paper5temporal} demonstrated that an offline-peek
grid-search procedure --- refitting HygienePrio's scorer weights on
calibration seeds at the same window being scored --- lifts
Precision@50 by up to $+21.3$~pp over the fixed Paper~4
weights~\cite{paper4hygieneprio} at the canonical $(K=50, \lambda=3)$
operating point. The procedure is not deployable: it requires
future-window data at scoring time. This paper asks whether any
deployable recalibration procedure can approach the offline ceiling
at moderate capacity while avoiding the absolute-floor collapse that
Paper~6~\cite{paper6capacity} documented at high capacity.

We evaluate six procedures in a pre-registered sequence, each
addressing the failure mode the previous one exposed:
\textbf{E1 lag-1 online} (rolling-history grid search at window $w-1$);
\textbf{E2 multi-window smoothing} (EWMA-3, trailing-mean-3 over
calibration history);
\textbf{E3 single-$\tau$ adaptive} (magnitude change-point detector
on the calibration-target shift);
\textbf{E4 threshold sensitivity sweep} ($\tau \in \{0.02,\ldots,0.10\}$);
\textbf{E5 capacity-aware $\tau$} (pre-registered $\tau_K$ vector);
\textbf{E6 one-sided CUSUM} (accumulator detector with pre-registered
slack $k$ and threshold $h$).
Each experiment is governed by its own pre-registration with locked
hypotheses, decision rules, and stop rules; 27 hypotheses across six
experiments are reported with honest outcome.

\textbf{Five substantive findings emerge.}
\textit{(1) E1 lag-1 recovers the offline ceiling at moderate capacity
($K \leq 100$: per-cell recovery ratio $\rho \approx 1.0$) but harms
performance at $K = 200$ ($\rho = -0.66$).}
\textit{(2) E2 smoothing amplifies the K=200 hazard rather than
reversing it (EWMA-3 $\rho = -0.94$); the naive "more history reduces
variance" prior is falsified.}
\textit{(3) E3 single-$\tau$ adaptive avoids the K=200 hazard
($\rho = +0.7$ pp) but pays a K=50 over-fire cost ($-5.3$~pp).
E4 shows no single $\tau$ in the swept grid satisfies both pre-registered
tolerances; the magnitude detector cannot reach the joint feasibility
region.}
\textit{(4) E5 capacity-aware $\tau$ reaches the feasibility region
but collapses to a static capacity-conditional gate: at K=200 the
cell mean equals the gate's exactly because the small $\tau_{200}$
fires every post-W1 window.}
\textit{(5) E6 CUSUM is the first procedure that beats the static
gate at K=200, by $+0.9$~pp, because the accumulator's slack filters
single-window noise excursions that the magnitude detector fires on.
This is the only operationally meaningful gain over the static gate in
the entire six-experiment sequence. CUSUM with a single $(k, h)$ pair
nevertheless over-fires at K=50 ($-3.9$~pp).}

\textbf{The deployable conclusion is conditional on capacity:}
at $K \leq 100$ deploy lag-1 directly; at $K = 200$ deploy CUSUM
$(k=0.04, h=0.10)$ for $+0.9$~pp over the static gate; for
deployments operating across both regimes a per-K $(k_K, h_K)$
vector is required. All results are bounded to the synthetic EEHDA
evaluation context; external validation on real fleet telemetry
with longitudinal exploitation outcome data is required before any
deployment recommendation. The frozen 19{,}500-row evaluation
corpus across the six experiments is released for independent
re-analysis.
